% 第一代完整逐步模拟 — 由第 9 章 \input

\section{模拟设定（全书统一）}

\begin{simbox}[玩具文件与参数]
\SimPara

\vspace{0.5em}
\textbf{12 字节文件：}\par
\SimFileTab

\vspace{0.5em}
\textbf{本例用到的 gear 表（节选）：}\par
\SimGearTab

{\footnotesize 注：数值专为教学设计，与 ebbackup 生产表不同，但算法步骤完全一致。}
\end{simbox}

\section{Chunk 1：从 pos=0 到第一次切点}

\subsection{Step A — 算 scan 区间}

\begin{simbox}[护栏模拟]
\begin{align*}
\texttt{pos} &= 0 \\
\texttt{scan\_start} &= 0 + \texttt{min} = 2 \\
\texttt{cut\_limit} &= \min(0+8,\ 12) = 8
\end{align*}
主循环 $i \in \{2,3,4,5,6,7\}$。
\end{simbox}

\subsection{Step B — 暖机（InitWindowHash）}

窗口 $w=2$，在 \texttt{scan\_start=2} 前窗口里是字节 \texttt{0x10, 0x20}：

\begin{simbox}[暖机算 $h$]
\begin{enumerate}[nosep]
  \item 读 \texttt{0x10}：$h = \mathrm{gear[0x10]} = 16$
  \item 读 \texttt{0x20}：$h = (16 \ll 1) + \mathrm{gear[0x20]} - \mathrm{gear[0x10]}
        = 32 + 2 - 16 = \mathbf{18}$
\end{enumerate}
暖机结束，进入 for 循环时 $h=18$。
\end{simbox}

\subsection{Step C — 逐字节 probe}

\begin{table}[htbp]
\centering
\caption{Chunk 1 逐步模拟（mask=3，看 $h \bmod 4$）}
\small
\begin{tabular}{@{}cccccp{3.8cm}@{}}
\toprule
\textbf{$i$} & \textbf{当前 $h$} & \textbf{$h\&3$} & \textbf{切？} & \textbf{读字节} & \textbf{滚动后 $h$} \\
\midrule
2 & 18 & 2 & 否 & \texttt{0x30} &
$(18\!\ll\!1)+0-16=20$ \\
3 & 20 & \textcolor{CutRed}{\textbf{0}} & \textcolor{CutRed}{\textbf{是}} & — & — \\
\bottomrule
\end{tabular}
\end{table}

\begin{figure}[htbp]
\centering
\begin{tikzpicture}[x=0.65cm,y=0.6cm]
  \foreach \i/\lab in {0/10,1/20,2/30,3/40,4/50,5/60,6/70,7/80,8/90,9/A0,10/B0,11/C0} {
    \node[byte, minimum width=0.6cm] at (\i,0) {\lab};
    \node[font=\tiny, below] at (\i,-0.45) {\i};
  }
  \fill[Gen1Blue!15] (0,-0.35) rectangle (2.65,0.65);
  \node[font=\scriptsize, Gen1Blue] at (1.3,0.95) {chunk 1（3 字节）};
  \draw[cutline] (3.35,-0.55) -- (3.35,0.75);
  \node[font=\scriptsize, CutRed, below] at (3.35,-0.75) {切点 $i=3$};
\end{tikzpicture}
\caption{Chunk 1 = 字节 0..2（切点 3 是下一块起点）。}
\end{figure}

\section{Chunk 2：从 pos=3 继续}

\subsection{Step D — 新区间与暖机}

\begin{simbox}[Chunk 2 起点]
\texttt{pos=3},\ \texttt{scan\_start=5},\ \texttt{cut\_limit=11}。\\
暖机窗口：字节 \texttt{0x40, 0x50} →
$h = 64$；$(64 \ll 1)+5-64 = \mathbf{69}$。
\end{simbox}

\subsection{Step E — 继续 probe}

\begin{table}[htbp]
\centering
\caption{Chunk 2 逐步模拟}
\small
\begin{tabular}{@{}cccccp{3.5cm}@{}}
\toprule
\textbf{$i$} & \textbf{$h$} & \textbf{$h\&3$} & \textbf{切？} & \textbf{读字节} & \textbf{滚动后 $h$} \\
\midrule
5 & 69 & 1 & 否 & \texttt{0x60} & $(69\!\ll\!1)+13-5=146$ \\
6 & 146 & 2 & 否 & \texttt{0x70} & $(146\!\ll\!1)+9-5=296$ \\
7 & 296 & \textcolor{CutRed}{\textbf{0}} & \textcolor{CutRed}{\textbf{是}} & — & — \\
\bottomrule
\end{tabular}
\end{table}

\section{Chunk 3 预览与 max-cut 演示}

若 Chunk 2 之后一直不命中，走到 \texttt{cut\_limit} 会 \textbf{max 强制切}。
本文件在 $i=7$ 已切，Chunk 3 从位置 7 开始，留给读者自练（习题见附录）。

\begin{simbox}[自练：Chunk 3 从 pos=7 开始]
\texttt{scan\_start=9}，暖机 $h=31$。\\
$i=9$：$31\&3=3$ 不切；滚动后 $h=48$。\\
$i=10$：$48\&3=0$ → \textbf{切点 = 10}，Chunk 3 长度 3。
\end{simbox}

\section{模拟小结}

\begin{genonebox}[第一代 9 步模拟链]
\begin{enumerate}[nosep]
  \item 定 pos / scan\_start / cut\_limit
  \item 暖机得初始 $h$
  \item 每个 $i$：先查 $(h \& mask)==0?$
  \item 未切：读 \texttt{文件[i]}，按 gear 表滚动
  \item 命中：切点 = $i$，输出 chunk
  \item 下一块：pos $\leftarrow$ 切点，重复
\end{enumerate}
\end{genonebox}
